\documentclass{article}
\usepackage{xfp}
\usepackage{tikz}
\usetikzlibrary{math}

\input{../../tikzlibrarymc.code.tex}

\begin{document}

% -----------------------------------------------
\tikzmath{
  print{\ \\===== Power Method =====};
  \a{1,1} = 4; \a{1,2} = -1; \a{1,3} = 1; 
  \a{2,1} = -1; \a{2,2} = 3; \a{2,3} = -2; 
  \a{3,1} = 1; \a{3,2} = -2; \a{3,3} = 3; 
}

\tikzset{
  % mat/identity={\a,3},
  mat/power method={\a,3,\b,\c},
  mat/show={\a,3,3},
  mat/show={\c,3,1},
}

\tikzmath{
  print{\ \\eigen value:\b};
  print{\ \\status:\status};
}

\end{document}